%%%%%%%%%%%%%%%%%%%%%%%%%%%%%%%%%%%%%%%%%%%%%%%%%%%%%%%%%%%%%%%%%%%%%%%%%%%%%%%
%
%  LaTeX-Vorlage der Hochschule Karlsruhe
%  fuer Skripte, Abschlussarbeiten, Projektarbeiten
%
%  ---------------------------------------------------------------------------
%  ALLES, WAS ANGEPASST WERDEN MUSS, STEHT IM ABSCHNITT "ZENTRALE ANGABEN".
%  Darunter ist normalerweise nichts zu aendern.
%  ---------------------------------------------------------------------------
%
%  Kompilieren:  latexmk -pdf document.tex
%  (ruft pdflatex, biber, makeglossaries und makeindex in richtiger
%   Reihenfolge und so oft wie noetig auf)
%
%  Hinweis zu Kapitelbildern: Bilder sollten ein Seitenverhaeltnis von
%  2:1 haben, z.B. 920 x 460 Pixel.
%
%  Basiert auf dem "Legrand Orange Book Template"
%  http://www.latextemplates.com/template/the-legrand-orange-book
%  von Mathias Legrand, modifiziert von Vel (vel@latextemplates.com).
%  Lizenz des Originals: CC BY-NC-SA 3.0
%  http://creativecommons.org/licenses/by-nc-sa/3.0/
%  ACHTUNG: Diese Vorlage ist eine Bearbeitung und steht daher ebenfalls
%  unter CC BY-NC-SA 3.0 -- inklusive der Non-Commercial-Klausel.
%
%%%%%%%%%%%%%%%%%%%%%%%%%%%%%%%%%%%%%%%%%%%%%%%%%%%%%%%%%%%%%%%%%%%%%%%%%%%%%%%

\documentclass[11pt,fleqn]{book} % Grundschriftgroesse, linksbuendige Formeln

\usepackage{xcolor} % muss vor der Farbwahl geladen sein


%%%%%%%%%%%%%%%%%%%%%%%%%%%%%%%%%%%%%%%%%%%%%%%%%%%%%%%%%%%%%%%%%%%%%%%%%%%%%%%
%  ZENTRALE ANGABEN -- HIER ANPASSEN
%%%%%%%%%%%%%%%%%%%%%%%%%%%%%%%%%%%%%%%%%%%%%%%%%%%%%%%%%%%%%%%%%%%%%%%%%%%%%%%

%------------------------------------------------------------------------------
%  1. Fakultaet: genau eine Zeile aktiv lassen
%------------------------------------------------------------------------------
% Die Hausfarben sind unten alle vordefiniert. Zum Wechseln der Fakultaet
% muss nur diese eine \colorlet-Zeile geaendert werden.

\definecolor{HKA@eit}{RGB}{ 40,135, 50} % Elektro- und Informationstechnik
\definecolor{HKA@ab} {RGB}{ 25,130,130} % Architektur und Bauwesen
\definecolor{HKA@mmt}{RGB}{ 40,105,175} % Maschinenbau und Mechatronik
\definecolor{HKA@w}  {RGB}{ 30, 70,150} % Wirtschaftswissenschaften
\definecolor{HKA@iwi}{RGB}{100, 55,140} % Informatik und Wirtschaftsinformatik
\definecolor{HKA@imm}{RGB}{140, 45,130} % Informationsmanagement und Medien

\colorlet{ThemeColor}{HKA@eit}          % <<<<< HIER Fakultaet waehlen

%------------------------------------------------------------------------------
%  2. Angaben zum Dokument
%------------------------------------------------------------------------------
\newcommand{\scriptTitle}   {Titel des Skripts}
\newcommand{\scriptSubject} {Thema des Skripts}
\newcommand{\scriptAuthor}  {Autorenname}
\newcommand{\scriptFaculty} {Elektro- und \\ Informationstechnik}
\newcommand{\scriptCourse}  {Vorlesung XYZ}
\newcommand{\scriptKeywords}{Keyword1, Keyword2, Keyword3}

%------------------------------------------------------------------------------
%  3. Angaben zur Veroeffentlichung
%------------------------------------------------------------------------------
\newcommand{\scriptPublisher}{NAME}
\newcommand{\scriptWebsite}  {https://book-website.com}
\newcommand{\scriptLicense}  {CC BY-NC-SA 3.0
                              (http://creativecommons.org/licenses/by-nc-sa/3.0/)}
\newcommand{\scriptEdition}  {Erste Ausgabe, Stand: \today}
\newcommand{\scriptYear}     {2026}

%------------------------------------------------------------------------------
%  4. Kapitelbilder
%------------------------------------------------------------------------------
% Bilder liegen in Resources/Pictures/ und werden ohne diesen Ordnernamen angegeben.
% Wer generell keine Kapitelbilder moechte, setzt hier \usechapterimagefalse
% (die Zeile also einfach einkommentieren). Einzelne Kapitel ohne Bild
% funktionieren ebenfalls -- dann bleibt der Kopfbereich schlicht weiss.

\newcommand{\tocChapterImage}    {chapter_image_1.jpg} % Inhaltsverzeichnis
\newcommand{\backChapterImage}   {chapter_image_2.jpg} % Verzeichnisse am Ende


%%%%%%%%%%%%%%%%%%%%%%%%%%%%%%%%%%%%%%%%%%%%%%%%%%%%%%%%%%%%%%%%%%%%%%%%%%%%%%%
%  AB HIER NORMALERWEISE NICHTS AENDERN
%%%%%%%%%%%%%%%%%%%%%%%%%%%%%%%%%%%%%%%%%%%%%%%%%%%%%%%%%%%%%%%%%%%%%%%%%%%%%%%

%------------------------------------------------------------------------------
%  Pakete, Layout, Makros
%------------------------------------------------------------------------------
% Reihenfolge beachten: ThemeColor ist oben bereits definiert, weil
% structure_document.tex die Farbe im Praeambel-Code verwendet.

\input{structure_document.tex} % Seitenlayout, Schriften, Ueberschriften, Boxen
%%%%%%%%%%%%%%%%%%%%%%%%%%%%%%%%%%%%%%%%%%%%%%%%%%%%%%%%%%%%%%%%%%%%%%%%%%%%%%%
%
%  customcommands.tex -- Zusatzpakete und eigene Befehle
%
%  Arbeitsteilung der Vorlage:
%    structure_document.tex  Seitenlayout, Schriften, Ueberschriften, Boxen
%    customcommands.tex      dieses File: Komfortpakete und eigene Makros
%    acronyms.tex            Glossar- und Abkuerzungseintraege
%
%  ACHTUNG: Diese Datei wird NACH structure_document.tex eingebunden und
%  setzt Definitionen von dort voraus -- insbesondere \avantfont (Schrift der
%  Titelseite) und die Farbe ThemeColor.
%
%%%%%%%%%%%%%%%%%%%%%%%%%%%%%%%%%%%%%%%%%%%%%%%%%%%%%%%%%%%%%%%%%%%%%%%%%%%%%%%


%==============================================================================
%  1. PAKETE
%==============================================================================

%--- Grafiken -----------------------------------------------------------------
\usepackage{graphicx}
\graphicspath{{../Pictures/}} % alle Bilder liegen hier; Angabe im Text ohne "Pictures/"
\usepackage{tikz}          % Titelseite, Kapitelkoepfe, Boxen
\usepackage{etoolbox}      % robuste Tests wie \ifdefempty (s. Titelseite)

%--- Listen und Tabellen ------------------------------------------------------
\usepackage{enumitem}      % Listen anpassbar
\setlist{nolistsep}        % kompaktere Abstaende in itemize/enumerate
\usepackage{booktabs}      % \toprule, \midrule, \bottomrule statt \hline

%--- Gleitobjekte -------------------------------------------------------------
% Zwingt Abbildungen und Tabellen, vor dem naechsten \section zu erscheinen.
% Wirkt nicht auf \subsection -- dort bei Bedarf \FloatBarrier von Hand setzen.
\usepackage[section]{placeins}

%--- Nur fuer die Beispielkapitel ---------------------------------------------
% \lipsum erzeugt Blindtext. Fuer die Abgabe-/Release-Version kann diese Zeile zusammen mit Document/99_example.tex entfernt werden.
\usepackage{lipsum}

%==============================================================================
%  2. ABBILDUNGEN
%==============================================================================

%------------------------------------------------------------------------------
%  \ofig -- "One FIGure": Abbildung mit einem Bild
%------------------------------------------------------------------------------
%  Nutzung:
%    \ofig{Bildunterschrift}{dateiname}{breite}
%    \ofig[eigenesLabel]{Bildunterschrift}{dateiname}{breite}
%
%  Verweis darauf:  \ref{fig:dateiname}  bzw.  \ref{fig:eigenesLabel}
%
%  Der Dateiname wird ohne Endung und ohne "Pictures/" angegeben.
%  Breite z.B. \linewidth, 0.6\linewidth, 8cm.
%
%  Ohne optionales Label wird der Dateiname als Label benutzt. Das ist bequem,
%  geht aber schief, sobald der Dateiname Unterstriche oder Punkte enthaelt --
%  dann das optionale Argument verwenden.
%------------------------------------------------------------------------------
\NewDocumentCommand{\ofig}{o m m m}{%
	\begin{figure}[htbp]
		\centering
		\includegraphics[width=#4]{#3}%
		\caption{#2}%
		\IfValueTF{#1}{\label{fig:#1}}{\label{fig:#3}}%
	\end{figure}%
}

%------------------------------------------------------------------------------
%  \pfig -- "Placeholder FIGure": Abbildung mit dem Platzhalterbild
%------------------------------------------------------------------------------
%  Nutzung:  \pfig{Bildunterschrift}
%  Praktisch, solange das echte Bild noch fehlt, damit Umbruch und
%  Nummerierung schon stimmen.
%------------------------------------------------------------------------------
\NewDocumentCommand{\pfig}{o m}{%
	\IfValueTF{#1}{\ofig[#1]{#2}{placeholder}{\linewidth}}%
	              {\ofig{#2}{placeholder}{\linewidth}}%
}


%==============================================================================
%  3. TITELSEITE
%==============================================================================

%------------------------------------------------------------------------------
%  Stellschrauben der Titelseite
%------------------------------------------------------------------------------
%  Diese drei Werte steuern das Layout. Die Bildhoehe wird daraus berechnet,
%  damit die Titelseite auch bei anderem Papierformat als A4 stimmt:
%  Die Oberkante des Bildes liegt bei 3/5 der Seitenhoehe, der farbige Balken
%  steht darunter am Seitenfuss.
%------------------------------------------------------------------------------
\newcommand{\hkaCoverImage}{cover_image.jpg} % leer lassen => kein Titelbild

\newlength{\hkaColorBlockHeight}
\setlength{\hkaColorBlockHeight}{3cm}        % Hoehe des farbigen Balkens unten

\newlength{\hkaCoverImageHeight}
\setlength{\hkaCoverImageHeight}{\dimexpr 0.6\paperheight - \hkaColorBlockHeight\relax}
% Bei A4: 0.6 * 29.7cm = 17.82cm, minus 3cm Balken = 14.82cm.

%------------------------------------------------------------------------------
%  \createTikZTitlePage -- baut die Titelseite auf
%------------------------------------------------------------------------------
%  Nutzung:
%    \createTikZTitlePage{Fakultaet}{Titel}{Autor(en)}
%
%  Der Fakultaetsname darf mit \\ umgebrochen werden.
%  Die Farbe kommt aus ThemeColor, wird also in document.tex festgelegt.
%------------------------------------------------------------------------------
\NewDocumentCommand{\createTikZTitlePage}{m m m}{%
	\begingroup
	\thispagestyle{empty} % keine Kopf-/Fusszeile auf der Titelseite
	\begin{tikzpicture}[remember picture,overlay]

	%--- Ankerpunkte (absolut auf der Seite) --------------------------------
		\coordinate (UoAS)         at ($ (current page.north west) + (2,-2) $);
		\coordinate (Faculty)      at ($ (UoAS) + (0,-2) $);
		\coordinate (Title)        at ($ (current page.north west) + (2,-11.5) $);
		\coordinate (Authors)      at ($ (current page.south west) + (2,1.5) $);
		\coordinate (BackgroundPic) at ($ (current page.south) + (0,\hkaColorBlockHeight) $);
		% Ankerpunkte des "+|KA"-Logos
		\coordinate (H) at ($ (current page.north east) + (8,-0.8) $);
		\coordinate (K) at ($ (H) + (0,-18) $);
		\coordinate (A) at ($ (K) + (0,-8.6) $);

	%--- Titelbild ----------------------------------------------------------
	% Ist \hkaCoverImage leer, bleibt die Flaeche frei -- so laesst sich die
	% Vorlage auch ohne mitgeliefertes Bild uebersetzen.
		\ifdefempty{\hkaCoverImage}{}{%
			\node[above, inner sep=0pt] (Hinterlegbild) at (BackgroundPic)
				{\includegraphics[height=\hkaCoverImageHeight]{\hkaCoverImage}};%
		}

	%--- Farbiger Balken am Seitenfuss --------------------------------------
		\path[fill=ThemeColor] (current page.south west)
			rectangle ($ (current page.south east) + (0,\hkaColorBlockHeight) $);

	%--- Textblöcke ---------------------------------------------------------
		% Hochschule
		\node[text width = 10cm, ThemeColor, below right] at (UoAS)
			{\avantfont{\textbf{Hochschule Karlsruhe}}\\
			 \avantfont{University of}\\
			 \avantfont{Applied Sciences}};
		% Fakultät
		\node[text width = 10cm, ThemeColor, below right] at (Faculty)
			{\avantfont{Fakultät für}\\ \avantfont{\textbf{#1}}};
		% Titel (und ggf. Untertitel)
		\node[text width = 10cm, ThemeColor, above right] at (Title) {#2};
		% Autor(en) -- steht im farbigen Balken, daher weiß
		\node[text width = 10cm, white, right] at (Authors) {{\large #3}};

	%--- Logo "+|KA" --------------------------------------------------------
	% Als TikZ-Pfade gezeichnet, damit das Logo scharf und in ThemeColor
	% bleibt. "+(x,y)" ist relativ zum zuletzt gesetzten Ankerpunkt, ohne
	% diesen zu verschieben.
		\path (H) [fill=ThemeColor, scale = 0.45, rounded corners = 1]
			+(0,-2.95) -- +(3,-2.95) -- +(3,0) -- +(4.2,0) -- +(4.2,-2.95)
			-- +(7.2,-2.95) -- +(7.2,-3.95) -- +(4.2,-3.95) -- +(4.2,-6.9)
			-- +(3,-6.9) -- +(3,-3.95) -- +(0,-3.95) -- cycle;
		\path (H) [fill=ThemeColor, scale = 0.45, rounded corners = 1]
			+(7.8,0) -- +(9,0) -- +(9,-6.9) -- +(7.8,-6.9) -- cycle;
		\path (K) [fill=ThemeColor, scale = 0.45, rounded corners = 1]
			+(3,0) -- +(4.2,0) -- +(4.2,-2.8) -- +(9,0) -- +(9,-1.3)
			-- +(5.2,-3.45) -- +(9,-5.6) -- +(9,-6.9) -- +(4.2,-4.1)
			-- +(4.2,-6.9) -- +(3,-6.9) -- cycle;
		\path (A) [fill=ThemeColor, scale = 0.45, rounded corners = 1]
			+(6,0) -- +(6.65,0) -- +(10.05,-6.9) -- +(8.75,-6.9) -- +(6,-1.3)
			-- +(4.8,-3.7) -- +(7.52,-3.7) -- +(8.85,-4.8) -- +(4.3,-4.8)
			-- +(3.25,-6.9) -- +(1.95,-6.9) -- +(5.35,0) -- cycle;

	\end{tikzpicture}%
	\vfill
	\endgroup
}     % eigene Befehle (\ofig, \pfig, Titelseite)
\input{acronyms.tex}           % Glossareintraege und Abkuerzungen

%------------------------------------------------------------------------------
%  PDF-Metadaten
%------------------------------------------------------------------------------
% Wichtig fuer Auffindbarkeit und Barrierefreiheit des fertigen PDFs.

\hypersetup{
	pdftitle    = {\scriptTitle},
	pdfauthor   = {\scriptAuthor},
	pdfsubject  = {\scriptSubject},
	pdfkeywords = {\scriptKeywords},
	pdflang     = {de-DE},
}

\date{\today}


\begin{document}

%------------------------------------------------------------------------------
%  VORSPANN -- roemische Seitenzahlen
%------------------------------------------------------------------------------
% \frontmatter sorgt fuer roemische Seitenzahlen (i, ii, iii ...) und dafuer,
% dass der Vorspann die Kapitelnummerierung nicht verbraucht.

\frontmatter

%--- Titelseite ---------------------------------------------------------------
\createTikZTitlePage{\scriptFaculty}
	{{\huge\textbf{\scriptTitle\\}}{Stand \today}}
	{\scriptAuthor}

%--- Impressum / Copyright ----------------------------------------------------
\newpage
\thispagestyle{empty}
~\vfill

\noindent Copyright \copyright\ \scriptYear\ \scriptAuthor

\noindent Herausgegeben von \scriptPublisher

\noindent \url{\scriptWebsite}

\medskip

\noindent Dieses Skript wird als Lehrmaterial für die Vorlesung \scriptCourse\
an der Hochschule Karlsruhe verwendet. Es wird unter der folgenden Lizenz
veröffentlicht: \scriptLicense.

\medskip

\noindent \textit{\scriptEdition}

%--- Inhaltsverzeichnis -------------------------------------------------------
\chapterimage{\tocChapterImage}

\pagestyle{empty} % keine Kopf-/Fusszeilen im Inhaltsverzeichnis

\tableofcontents

% Abbildungs- und Tabellenverzeichnis -- bei Bedarf einkommentieren:
%\cleardoublepage
%\addcontentsline{toc}{chapter}{Abbildungsverzeichnis}
%\listoffigures
%
%\cleardoublepage
%\addcontentsline{toc}{chapter}{Tabellenverzeichnis}
%\listoftables

\cleardoublepage % erzwingt Start des ersten Kapitels auf einer rechten Seite

%------------------------------------------------------------------------------
%  HAUPTTEIL -- arabische Seitenzahlen, nummerierte Kapitel
%------------------------------------------------------------------------------
\mainmatter

\pagestyle{fancy} % Kopf-/Fusszeilen wieder aktivieren

%--- Teil 1 -------------------------------------------------------------------
\part{TEIL 1}

%%%%%%%%%%%%%%%%%%%%%%%%%%%%%%%%%%%%%%%%%%%%%%%%%%%%%%%%%%%%%%%%%%%%%%%%%%%%%%%
%  Kapitel 1 -- Einleitung
%
%  Dieses Kapitel ist als Geruest gedacht: Die Zwischenueberschriften decken
%  ab, was eine Einleitung ueblicherweise leistet. Text ersetzen, nicht
%  benoetigte Abschnitte loeschen.
%
%  Bilder liegen in Pictures/ und werden OHNE diesen Ordnernamen und OHNE
%  Dateiendung angegeben -- \graphicspath erledigt den Rest.
%%%%%%%%%%%%%%%%%%%%%%%%%%%%%%%%%%%%%%%%%%%%%%%%%%%%%%%%%%%%%%%%%%%%%%%%%%%%%%%

\chapterimage{chapter_image_1.jpg} % Bild im Kapitelkopf (Seitenverhaeltnis 2:1)

\chapter{Einleitung}
\label{ch:einleitung}
\index{Einleitung}


%==============================================================================
\section{Motivation}
\label{sec:motivation}
%==============================================================================

Hier steht, warum das Thema überhaupt behandelt wird: das Problem, die
Ausgangslage, der Anlass. Anführungszeichen setzt man mit \enquote{enquote},
dann stimmen sie automatisch in jeder Sprache -- verschachtelt funktioniert
das ebenfalls, etwa \enquote{ein Zitat mit \enquote{Zitat darin}}.

Zahlen mit Einheiten gehören in \verb|\SI|, damit Abstand und Dezimalzeichen
einheitlich bleiben: Die Messung ergab \SI{12345.67}{\metre\per\second} bei
einer Spannung von \SI{3.3}{\volt}. Im Quelltext schreibt man dabei immer
einen Punkt -- die Ausgabe als Komma erledigt \texttt{siunitx}.

\lipsum[1] % Blindtext -- vor der Abgabe entfernen


%==============================================================================
\section{Zielsetzung}
\label{sec:zielsetzung}
%==============================================================================

Was soll am Ende erreicht sein? Idealerweise so konkret, dass sich im Fazit
überprüfen lässt, ob es erreicht wurde.

Auf andere Stellen verweist man mit \verb|\cref|. Das Wort davor liefert der
Befehl bereits mit -- man schreibt also nur \verb|\cref{ch:einleitung}| und
erhält \cref{ch:einleitung}. Ein zusätzliches \enquote{Kapitel} davor wäre
doppelt. Für Abschnitte, Abbildungen und Tabellen gilt dasselbe:
\cref{sec:motivation} und \cref{fig:beispielbild}.

\ofig[beispielbild]{Ein Platzhalterbild. Die Bildunterschrift steht bei
Abbildungen unter dem Bild und endet ohne Punkt}{placeholder}{0.6\linewidth}

Abkürzungen werden bei der ersten Verwendung automatisch ausgeschrieben:
\acrfull{gcd}. Ab dann genügt die Kurzform \acrshort{gcd}. Das
Abkürzungsverzeichnis am Ende des Dokuments füllt sich von selbst -- neue
Einträge kommen in \texttt{acronyms.tex}.


%==============================================================================
\section{Aufbau der Arbeit}
\label{sec:aufbau}
%==============================================================================

Ein kurzer Wegweiser durch die folgenden Kapitel. Quellen werden mit
\verb|\cite| zitiert \cite{book_key}; mit Seitenangabe sieht das so
aus \cite[162]{book_key}. Die Einträge dafür stehen in der Datei
\texttt{bibliography.bib}.

\lipsum[2] % Blindtext -- vor der Abgabe entfernen


%--- Teil 2: Formatbeispiele --------------------------------------------------
% In der Abgabe-/Release-Version loeschen.
\part{TEIL 2}

%%%%%%%%%%%%%%%%%%%%%%%%%%%%%%%%%%%%%%%%%%%%%%%%%%%%%%%%%%%%%%%%%%%%%%%%%%%%%%%
%  Formatbeispiele -- Schaufenster aller Umgebungen der Vorlage
%
%  Dieses Kapitel dient als Nachschlagewerk beim Schreiben und gehoert NICHT
%  in die Abgabe. Vor der finalen Version einfach den zugehoerigen \part und
%  \input in document.tex loeschen.
%
%  Die Namen der Umgebungen sind englisch (\begin{theorem}), die Ausgabe ist
%  deutsch (\enquote{Satz}) -- festgelegt in structure_document.tex.
%%%%%%%%%%%%%%%%%%%%%%%%%%%%%%%%%%%%%%%%%%%%%%%%%%%%%%%%%%%%%%%%%%%%%%%%%%%%%%%

\chapter{Formatierungen und Vorlagen}
\label{ch:formatbeispiele}


%==============================================================================
\section{Tabellen}
\index{Tabelle}
%==============================================================================

Tabellen bekommen ihre Beschriftung \emph{über} den Inhalt. Senkrechte Linien
sind unüblich; \texttt{booktabs} liefert mit \verb|\toprule|, \verb|\midrule|
und \verb|\bottomrule| die passenden waagerechten.

Für Zahlenspalten lohnt sich der Spaltentyp \texttt{S} aus \texttt{siunitx}:
Er richtet am Dezimalzeichen aus und setzt automatisch das deutsche Komma.
Standardmäßig gruppiert \texttt{siunitx} auch die Nachkommastellen in
Dreierblöcken (\enquote{0,000\,326\,2}). Wer das nicht möchte, ergänzt
\texttt{group-digits=integer} -- entweder je Spalte wie unten oder global
über \verb|\sisetup| in \texttt{structure\_document.tex}.

\begin{table}[htbp]
	\centering
	\caption{Tabellenbeschriftung -- steht oben und endet ohne Punkt}
	\label{tab:beispiel}
	\begin{tabular}{l S[table-format=1.7, group-digits=integer] S[table-format=1.3]}
		\toprule
		\textbf{Behandlung} & {\textbf{Messwert 1}} & {\textbf{Messwert 2}} \\
		\midrule
		Behandlung 1 & 0.0003262 & 0.562 \\
		Behandlung 2 & 0.0015681 & 0.910 \\
		Behandlung 3 & 0.0009271 & 0.296 \\
		\bottomrule
	\end{tabular}
\end{table}

Verwiesen wird mit \verb|\cref|: \cref{tab:beispiel} zeigt das Ergebnis.


%==============================================================================
\section{Abbildungen}
\index{Abbildung}
%==============================================================================

Abbildungen bekommen ihre Beschriftung \emph{unter} das Bild. Statt einer
kompletten \texttt{figure}-Umgebung genügt der Kurzbefehl \verb|\ofig| aus
\texttt{customcommands.tex}:

\begin{verbatim}
\ofig[label]{Bildunterschrift}{dateiname}{breite}
\end{verbatim}

Der Dateiname steht ohne Ordner und ohne Endung. Ohne das optionale Label wird
der Dateiname als Label benutzt -- was schiefgeht, sobald zwei Abbildungen
dasselbe Bild verwenden. Deshalb hier ein eigenes Label.

\ofig[beispielabbildung]{Abbildungsbeschriftung -- steht unten und endet ohne
Punkt}{placeholder}{0.5\linewidth}

%\cref{fig:beispielabbildung} wird damit automatisch nummeriert und verlinkt.

%==============================================================================
\section{Schaltpläne}
\index{Schaltplan}
%==============================================================================

Über \texttt{circuitikz} lassen sich Schaltpläne direkt im Quelltext
zeichnen -- sie bleiben dadurch beliebig skalierbar und passen zur Schrift
des Dokuments.

\begin{figure}[htbp]
	\centering
	\begin{circuitikz}[scale=0.9, transform shape]
		\draw (0,0) to[V, l=$U_\mathrm{e}$] (0,4);
		\draw (0,4) to[R, l=$R_1$, *-*] (4,4);
		\draw (4,4) to[R, l=$R_2$] (4,0);
		\draw (4,0) -- (0,0);
		\draw (4,4) to[short, -o] (5.5,4) node[right] {$U_\mathrm{a}$};
		\draw (4,0) to[short, -o] (5.5,0);
		\draw (0,0) node[ground]{};
	\end{circuitikz}
	\caption{Unbelasteter Spannungsteiler, gezeichnet mit \texttt{circuitikz}}
	\label{fig:spannungsteiler}
\end{figure}

Für \cref{fig:spannungsteiler} gilt
$U_\mathrm{a} = U_\mathrm{e} \cdot \frac{R_2}{R_1 + R_2}$.


%==============================================================================
\section{Quellenangaben}
\index{Zitieren}
%==============================================================================

Diese Aussage benötigt eine Quelle \cite{article_key}; diese hier ist
konkreter und nennt die Seite \cite[162]{book_key}.


%==============================================================================
\section{Listen}
\index{Liste}
%==============================================================================

Listen ordnen Informationen knapp und übersichtlich\footnote{So sieht eine
Fußnote aus.}.

\subsection{Nummerierte Liste}
\index{Liste!nummeriert}

\begin{enumerate}
	\item Erster Punkt
	\item Zweiter Punkt
	\item Dritter Punkt
\end{enumerate}

\subsection{Aufzählung}
\index{Liste!Aufzählung}

\begin{itemize}
	\item Erster Punkt
	\item Zweiter Punkt
	\item Dritter Punkt
\end{itemize}

\subsection{Beschreibungen}
\index{Liste!Beschreibung}

\begin{description}
	\item[Begriff] Erläuterung dazu
	\item[Wort] Definition dazu
	\item[Anmerkung] Ergänzung dazu
\end{description}


%==============================================================================
\section{Sätze}
\index{Satz}
%==============================================================================

% HINWEIS: Diese Section fehlte urspruenglich. Die beiden folgenden
% Unterabschnitte hingen dadurch faelschlich unter \section{Listen}, obwohl
% die Index-Eintraege bereits eine eigene Section voraussetzten.

\subsection{Mehrere Gleichungen}
\index{Satz!mehrere Gleichungen}

\begin{theorem}[Name des Satzes]
	In $E=\mathbb{R}^n$ sind alle Normen äquivalent. Es gilt:
	\begin{align}
		& \big| \|\mathbf{x}\| - \|\mathbf{y}\| \big| \leq \|\mathbf{x}-\mathbf{y}\| \\
		& \Big\|\sum_{i=1}^n \mathbf{x}_i\Big\| \leq \sum_{i=1}^n \|\mathbf{x}_i\|
		  \quad\text{mit $n$ endlich}
	\end{align}
\end{theorem}

\subsection{Einzeilig}
\index{Satz!einzeilig}

\begin{theorem}
	Eine Menge $\mathcal{D}(G)$ liegt dicht in $L^2(G)$, $|\cdot|_0$.
\end{theorem}


%==============================================================================
\section{Definitionen}
\index{Definition}
%==============================================================================

\begin{definition}[Name der Definition]
	Sei $E$ ein Vektorraum. Eine Norm auf $E$ ist eine Abbildung
	$\|\cdot\|$ von $E$ nach $\mathbb{R}^+=[0,+\infty[$ mit:
	\begin{align}
		& \|\mathbf{x}\|=0 \ \Rightarrow\ \mathbf{x}=\mathbf{0} \\
		& \|\lambda \mathbf{x}\| = |\lambda|\cdot\|\mathbf{x}\| \\
		& \|\mathbf{x}+\mathbf{y}\| \leq \|\mathbf{x}\|+\|\mathbf{y}\|
	\end{align}
\end{definition}


%==============================================================================
\section{Notationen}
\index{Notation}
%==============================================================================

\begin{notation}
	Sei $G$ eine offene Teilmenge von $\mathbb{R}^n$. Die Menge der
	Funktionen $\varphi$ mit
	\begin{enumerate}
		\item beschränktem Träger $G$ und
		\item beliebig oft differenzierbar
	\end{enumerate}
	bildet einen Vektorraum, bezeichnet mit $\mathcal{D}(G)$.
\end{notation}


%==============================================================================
\section{Bemerkungen}
\index{Bemerkung}
%==============================================================================

\begin{remark}
	Die hier verwendeten Begriffe sind in der Mathematik gebräuchlich.
	Vektorräume werden über dem Körper $\mathbb{K}=\mathbb{R}$ betrachtet;
	die Eigenschaften lassen sich auf $\mathbb{K}=\mathbb{C}$ übertragen.
\end{remark}


%==============================================================================
\section{Folgerungen}
\index{Folgerung}
%==============================================================================

\begin{corollary}[Name der Folgerung]
	Die hier verwendeten Begriffe sind in der Mathematik gebräuchlich.
	Vektorräume werden über dem Körper $\mathbb{K}=\mathbb{R}$ betrachtet.
\end{corollary}


%==============================================================================
\section{Thesen}
\index{These}
%==============================================================================

\subsection{Mehrere Gleichungen}
\index{These!mehrere Gleichungen}

\begin{proposition}[Name der These]
	Es gilt:
	\begin{align}
		& \big| \|\mathbf{x}\| - \|\mathbf{y}\| \big| \leq \|\mathbf{x}-\mathbf{y}\| \\
		& \Big\|\sum_{i=1}^n \mathbf{x}_i\Big\| \leq \sum_{i=1}^n \|\mathbf{x}_i\|
		  \quad\text{mit $n$ endlich}
	\end{align}
\end{proposition}

\subsection{Einzeilig}
\index{These!einzeilig}

\begin{proposition}
	Seien $f,g\in L^2(G)$. Gilt $(f,\varphi)_0=(g,\varphi)_0$ für alle
	$\varphi\in\mathcal{D}(G)$, so folgt $f=g$.
\end{proposition}


%==============================================================================
\section{Beispiele}
\index{Beispiel}
%==============================================================================

\subsection{Gleichung mit Text}
\index{Beispiel!Gleichung mit Text}

\begin{example}
	Sei $G=\{x\in\mathbb{R}^2:|x|<3\}$ und $x^0=(1,1)$. Betrachte
	\begin{equation}
		f(x)=\left\{\begin{aligned}
			& \mathrm{e}^{|x|} & & \text{falls $|x-x^0|\leq 1/2$}\\
			& 0                & & \text{falls $|x-x^0| > 1/2$}
		\end{aligned}\right.
	\end{equation}
	Die Funktion $f$ hat beschränkten Träger; man kann
	$A=\{x\in\mathbb{R}^2:|x-x^0|\leq 1/2+\epsilon\}$ wählen für alle
	$\epsilon\in\intoo{0}{5/2-\sqrt{2}}$.
\end{example}

\subsection{Reiner Textabsatz}
\index{Beispiel!Textabsatz}

\begin{example}[Name des Beispiels]
	\lipsum[2]
\end{example}


%==============================================================================
\section{Übungen}
\index{Übung}
%==============================================================================

\begin{exercise}
	Hier lässt sich eine Frage stellen, mit der Studierende ihren Lernstand
	überprüfen oder das Gelesene vertiefen können.
\end{exercise}


%==============================================================================
\section{Probleme}
\index{Problem}
%==============================================================================

\begin{problem}
	Wie hoch ist die durchschnittliche Fluggeschwindigkeit einer unbeladenen
	Schwalbe?
\end{problem}


%==============================================================================
\section{Wörterverzeichnis}
\index{Wörterverzeichnis}
%==============================================================================

\begin{vocabulary}[Wort]
	Definition des Wortes.
\end{vocabulary}


%==============================================================================
\section{Glossar und Abkürzungen}
\index{Glossar}
%==============================================================================

Einträge aus \texttt{acronyms.tex} werden über \verb|\gls| und \verb|\acr...|
eingebunden und landen automatisch in den Verzeichnissen am Dokumentende.

Die Auszeichnungssprache \Gls{latexg} eignet sich besonders für Dokumente
mit \gls{mathsg}.

Für eine Menge von Zahlen gibt es elementare Verfahren zur Bestimmung des
\acrlong{gcd}, abgekürzt \acrshort{gcd}. Ähnlich verfährt man beim
\acrfull{lcm}.


%------------------------------------------------------------------------------
%  ANHANG UND VERZEICHNISSE
%------------------------------------------------------------------------------
% \backmatter schaltet die Kapitelnummerierung ab, behaelt aber die
% arabische Seitenzaehlung bei -- genau das Richtige fuer die Verzeichnisse.

\backmatter

\chapterimage{\backChapterImage}

%--- Literaturverzeichnis -----------------------------------------------------
% Die Kopfzeile muss vor JEDEM Verzeichnis neu gesetzt werden. Wird das
% vergessen, tragen alle folgenden Seiten weiterhin die Kopfzeile des
% vorherigen Verzeichnisses.

\cleardoublepage
\phantomsection
\fancyhead[LE,LO]{\sffamily\normalsize\thepage}
\fancyhead[RE,RO]{\sffamily\normalsize\textbf{\bibname}}
% Kein \textcolor im Eintragstext: die Faerbung uebernimmt das
% numberless-Format von \titlecontents in structure_document.tex --
% im Inhaltsverzeichnis ThemeColor, auf den TEIL-Seiten weiss.
\addcontentsline{toc}{chapter}{Literaturverzeichnis}
\printbibliography[title=Literaturverzeichnis]

%--- Abkuerzungsverzeichnis ---------------------------------------------------
\cleardoublepage
\phantomsection
\fancyhead[RE,RO]{\sffamily\normalsize\textbf{Abkürzungsverzeichnis}}
\addcontentsline{toc}{chapter}{Abkürzungsverzeichnis}
\printglossary[type=\acronymtype,title={Abkürzungsverzeichnis}]

%--- Glossar ------------------------------------------------------------------
\cleardoublepage
\phantomsection
\fancyhead[RE,RO]{\sffamily\normalsize\textbf{Glossar}}
\addcontentsline{toc}{chapter}{Glossar}
\printglossary[title=Glossar]

%--- Index --------------------------------------------------------------------
\cleardoublepage
\phantomsection
\setlength{\columnsep}{0.75cm} % Abstand zwischen den beiden Index-Spalten
\fancyhead[RE,RO]{\sffamily\normalsize\textbf{\indexname}}
\addcontentsline{toc}{chapter}{Index}
\printindex

%------------------------------------------------------------------------------

\end{document}
